\documentclass[10pt, letterpaper]{article}

% Packages:
\usepackage[
    ignoreheadfoot, % set margins without considering header and footer
    top=2 cm, % seperation between body and page edge from the top
    bottom=2 cm, % seperation between body and page edge from the bottom
    left=2 cm, % seperation between body and page edge from the left
    right=2 cm, % seperation between body and page edge from the right
    footskip=1.0 cm, % seperation between body and footer
    % showframe % for debugging 
]{geometry} % for adjusting page geometry
\usepackage[explicit]{titlesec} % for customizing section titles
\usepackage{tabularx} % for making tables with fixed width columns
\usepackage{array} % tabularx requires this
\usepackage[dvipsnames]{xcolor} % for coloring text
\definecolor{primaryColor}{RGB}{0, 79, 144} % define primary color
\usepackage{enumitem} % for customizing lists
\usepackage{fontawesome5} % for using icons
\usepackage{amsmath} % for math
\usepackage[
    pdftitle={John Doe's CV},
    pdfauthor={John Doe},
    pdfcreator={LaTeX with RenderCV},
    colorlinks=true,
    urlcolor=primaryColor
]{hyperref} % for links, metadata and bookmarks
\usepackage[pscoord]{eso-pic} % for floating text on the page
\usepackage{calc} % for calculating lengths
\usepackage{bookmark} % for bookmarks
\usepackage{lastpage} % for getting the total number of pages
\usepackage{changepage} % for one column entries (adjustwidth environment)
\usepackage{paracol} % for two and three column entries
\usepackage{ifthen} % for conditional statements
\usepackage{needspace} % for avoiding page brake right after the section title
\usepackage{iftex} % check if engine is pdflatex, xetex or luatex

% Ensure that generate pdf is machine readable/ATS parsable:
\ifPDFTeX
    \input{glyphtounicode}
    \pdfgentounicode=1
    \usepackage[T1]{fontenc}
    \usepackage[utf8]{inputenc}
    \usepackage{lmodern}
\fi

\usepackage[default, type1]{sourcesanspro} 

% Some settings:
\AtBeginEnvironment{adjustwidth}{\partopsep0pt} % remove space before adjustwidth environment
\pagestyle{empty} % no header or footer
\setcounter{secnumdepth}{0} % no section numbering
\setlength{\parindent}{0pt} % no indentation
\setlength{\topskip}{0pt} % no top skip
\setlength{\columnsep}{0.15cm} % set column seperation
\makeatletter
% \let\ps@customFooterStyle\ps@plain % Copy the plain style to customFooterStyle
% \patchcmd{\ps@customFooterStyle}{\thepage}{
%     \color{gray}\textit{\small John Doe - Page \thepage{} of \pageref*{LastPage}}
% }{}{} % replace number by desired string
% \makeatother
% \pagestyle{customFooterStyle}

\titleformat{\section}{
    % avoid page braking right after the section title
    \needspace{4\baselineskip}
    % make the font size of the section title large and color it with the primary color
    \Large\color{primaryColor}
}{
}{
}{
    % print bold title, give 0.15 cm space and draw a line of 0.8 pt thickness
    % from the end of the title to the end of the body
    \textbf{#1}\hspace{0.15cm}\titlerule[0.8pt]\hspace{-0.1cm}
}[] % section title formatting

\titlespacing{\section}{
    % left space:
    -1pt
}{
    % top space:
    0.3 cm
}{
    % bottom space:
    0.2 cm
} % section title spacing

% \renewcommand\labelitemi{$\vcenter{\hbox{\small$\bullet$}}$} % custom bullet points
\newenvironment{highlights}{
    \begin{itemize}[
        topsep=0.10 cm,
        parsep=0.10 cm,
        partopsep=0pt,
        itemsep=0pt,
        leftmargin=0.4 cm + 10pt
    ]
}{
    \end{itemize}
} % new environment for highlights

\newenvironment{highlightsforbulletentries}{
    \begin{itemize}[
        topsep=0.10 cm,
        parsep=0.10 cm,
        partopsep=0pt,
        itemsep=0pt,
        leftmargin=10pt
    ]
}{
    \end{itemize}
} % new environment for highlights for bullet entries


\newenvironment{onecolentry}{
    \begin{adjustwidth}{
        0.2 cm + 0.00001 cm
    }{
        0.2 cm + 0.00001 cm
    }
}{
    \end{adjustwidth}
} % new environment for one column entries

\newenvironment{twocolentry}[2][]{
    \onecolentry
    \setcolumnwidth{\fill, 4.5cm}
    \begin{paracol}{2}
        % Left column: content starts immediately
        #1 \switchcolumn
        % Right column: date aligned at top
        \raggedleft #2
    \end{paracol}
    \endonecolentry
}{} % new environment for two column entries

\newenvironment{threecolentry}[3][]{
    \onecolentry
    \def\thirdColumn{#3}
    \setcolumnwidth{1 cm, \fill, 4.5 cm}
    \begin{paracol}{3}
    {\raggedright #2} \switchcolumn
}{
    \switchcolumn \raggedleft \thirdColumn
    \end{paracol}
    \endonecolentry
} % new environment for three column entries

\newenvironment{header}{
    \setlength{\topsep}{0pt}\par\kern\topsep\centering\color{primaryColor}\linespread{1.5}
}{
    \par\kern\topsep
} % new environment for the header

% \newcommand{\placelastupdatedtext}{% \placetextbox{<horizontal pos>}{<vertical pos>}{<stuff>}
%   \AddToShipoutPictureFG*{% Add <stuff> to current page foreground
%     \put(
%         \LenToUnit{\paperwidth-2 cm-0.2 cm+0.05cm},
%         \LenToUnit{\paperheight-1.0 cm}
%     ){\vtop{{\null}\makebox[0pt][c]{
%         \small\color{gray}\textit{Last updated in September 2024}\hspace{\widthof{Last updated in September 2024}}
%     }}}%
%   }%
% }%

% save the original href command in a new command:
\let\hrefWithoutArrow\href

% new command for external links:
\renewcommand{\href}[2]{\hrefWithoutArrow{#1}{\ifthenelse{\equal{#2}{}}{ }{#2 }\raisebox{.15ex}{\footnotesize \faExternalLink*}}}


% Education heading
\newcommand{\eduHeading}[4]{
  \item
  \begin{tabular*}{0.97\textwidth}[t]{l@{\extracolsep{\fill}}r}
    \textbf{#1} & #2 \\
    \textit{\small#3} & \textit{\small #4} \\
  \end{tabular*}
}

% Optional bullet list under each eduHeading
\newcommand{\eduItemListStart}{
  \begin{itemize}[leftmargin=0.35in, itemsep=0pt, topsep=2pt]
}
\newcommand{\eduItemListEnd}{\end{itemize}}

% Project entry environment
\newenvironment{projectentry}[3]{%
  \onecolentry
  % First row: project title (left) and date (right)
  \begin{tabular*}{\textwidth}{@{\extracolsep{\fill}} l r}
    \textbf{#1} & #2 \\
  \end{tabular*}
  % Advisor line (flush left, tight to previous line)
  \textit{#3} \par
  \vspace{0.2em} % small controlled spacing before highlights
}{%
  \endonecolentry
}

% Publication entry: left = title/authors/conference, right = date and location
\newcommand{\publicationEntry}[6]{%
  \begin{samepage}%
    \noindent
    \begin{tabular*}{\textwidth}{@{\extracolsep{\fill}} l r}
      % Left column: title + conference + authors
      \begin{minipage}[t]{0.68\textwidth}%
        \hypersetup{urlcolor=black}%
        \textbf{\href{#2}{#1}}\\[0.05cm]%
        \hypersetup{urlcolor=primaryColor}%
        #4\\[0.05cm]%
        #3%
      \end{minipage}%
      &
      % Right column: date + location stacked
      \begin{minipage}[t]{0.3\textwidth}%
        \raggedleft
        #5\\%
        #6%
      \end{minipage}%
    \end{tabular*}%
  \end{samepage}%
  \vspace{0.35cm}%
}

% Equal contribution note
\newcommand{\equalAuthor}[1]{#1\textsuperscript{*}}

% Footnote for equal contribution (place once at end of document)
\newcommand{\equalContributionNote}{%
  \vspace{0.5cm}
  {\footnotesize \textsuperscript{*} Authors contributed equally to this work.}
}

% Awards entry environment (using onecolentry)
\newcommand{\awardEntry}[3]{%
  \begin{samepage}%
    \noindent
    \begin{tabular*}{\textwidth}{@{\extracolsep{\fill}} l r}
      \textbf{#1} & \raggedleft #3 \\
    \end{tabular*}%
    \vspace{0.05cm}\\
    #2
  \end{samepage}%
  \vspace{0.35cm}%
}

\newlength{\skilllabelwidth}
\settowidth{\skilllabelwidth}{\textbf{Programming Languages}\quad}

\newcommand{\skillrow}[2]{%
  \par\noindent
  \hangindent=\skilllabelwidth
  \hangafter=1
  \makebox[\skilllabelwidth][l]{\textbf{#1}}#2\par
}

% Experience Entry
\newcommand{\experienceEntry}[5]{%
  \begin{samepage}%
    \noindent
    \begin{tabular*}{\textwidth}{@{\extracolsep{\fill}} l r}
      % Left column: title + organization
      \begin{minipage}[t]{0.68\textwidth}%
        \textbf{#1} \\[0.05cm]%
        \textit{#2}%
      \end{minipage}%
      &
      % Right column: date (top) and location (bottom)
      \begin{minipage}[t]{0.3\textwidth}%
        \raggedleft
        #4 \\[0.05cm]
        \textit{#3}%
      \end{minipage}%
    \end{tabular*}%
    % Reduce space before highlights
    \vspace{0.1cm}%
    % Optional bullet points
    #5
  \end{samepage}%
  \vspace{0.25cm}%
}

\begin{document}
    \newcommand{\AND}{\unskip
        \cleaders\copy\ANDbox\hskip\wd\ANDbox
        \ignorespaces
    }
    \newsavebox\ANDbox
    \sbox\ANDbox{}

    % \placelastupdatedtext
    \begin{header}
        \fontsize{30 pt}{30 pt}
        \textbf{Pushkal Mishra}

        \vspace{0.3 cm}

        \normalsize
        % \mbox{{\footnotesize\faMapMarker*}\hspace*{0.13cm}Your Location}%
        % \kern 0.25 cm%
        % \AND%
        % \kern 0.25 cm%
        \mbox{\hrefWithoutArrow{mailto:pumishra@ucsd.edu}{{\footnotesize\faEnvelope[regular]}\hspace*{0.13cm}pumishra@ucsd.edu}}%
        \kern 0.25 cm%
        \AND%
        % \kern 0.25 cm%
        % \mbox{\hrefWithoutArrow{tel:+1-858-373-7187}{{\footnotesize\faPhone*}\hspace*{0.13cm}+1-858-373-7187}}%
        % \kern 0.25 cm%
        % \AND%
        \kern 0.25 cm%
        \mbox{\hrefWithoutArrow{https://pushkalm11.github.io/}{{\footnotesize\faLink}\hspace*{0.13cm}pushkalm11.github.io}}%
        \kern 0.25 cm%
        \AND%
        \kern 0.25 cm%
        \mbox{\hrefWithoutArrow{https://www.linkedin.com/in/pushkal-mishra-014829218/}{{\footnotesize\faLinkedinIn}\hspace*{0.13cm}Pushkal Mishra}}%
        \kern 0.25 cm%
        \AND%
        \kern 0.25 cm%
        \mbox{\hrefWithoutArrow{https://github.com/PushkalM11}{{\footnotesize\faGithub}\hspace*{0.13cm}PushkalM11}}%
    \end{header}

    \vspace{0.3 cm - 0.3 cm}

\section{Biography}

\begin{onecolentry}
    I am a researcher focused on developing scalable deep learning systems that perceive and reason about the physical world. My work spans multi-modal radar perception, autonomous driving, and sensor fusion, with an emphasis on building frameworks that unify simulation, real-world data, and large language models for robust scene understanding and decision-making. I am particularly interested in creating modular architectures that enable plug-and-play integration of diverse sensors while improving interpretability and generalization under challenging environmental conditions.
\end{onecolentry}

\section{Education}
  \begin{itemize}[leftmargin=0.15in, label={}]

    % PhD
    \eduHeading
        {University of California San Diego}{San Diego, CA, USA}
        {PhD in Electrical and Computer Engineering}{Aug 2024 -- Present}
    \eduItemListStart
        \item Advisor: Prof. Dinesh Bharadia, \href{https://wcsng.ucsd.edu/}{WCSNG Lab}
        \item {
            Research Topics: \\
            Radar Sensing, Multimodal Representation Alignment, Autonomous Driving
        }
    \eduItemListEnd

    % B.Tech
    \eduHeading
        {Indian Institute of Technology Hyderabad}{Hyderabad, India}
        {B.Tech in Electrical Engineering, Minor in Computer Science}{Nov 2020 -- May 2024}
    \eduItemListStart
        \item PIs: Prof. Aditya Siripuram, Prof. Ayon Borthakur
        \item Graduated with CGPA: \textbf{9.33 / 10}
    \eduItemListEnd

    % High School
    \eduHeading
        {Royale Concorde International School}{Bangalore, India}
        {High School and Senior Secondary (CBSE)}{Jun 2005 -- Aug 2020}
    \eduItemListStart
        \item Grade 12: \textbf{96.2\%}, Grade 10: \textbf{90\%}
    \eduItemListEnd

  \end{itemize}


\section{Research Projects}

\begin{projectentry}{Radar Foundation Model}{Jan 2024 -- Present}{With Prof. Dinesh Bharadia at UCSD}
    \begin{highlights}
        \item Built a radar-to-text foundation model for captioning, sim-to-real transfer, and interpretability, unifying radar perception and downstream tasks for autonomous driving.
        \item This framework enabled LLMs to reason over structured driving scene descriptions.
        \item Devised novel evaluation metrics tailored for generative radar models.
        \item Large-scale pre-training using Kubernetes, signal processing on radars, and multi-modal ML pipelines to improve generalization under noise and adverse conditions.
        \item \textbf{Skills:} Multi-modal ML pipelines, radar signal processing, large-scale pre-training, generative model evaluation.
    \end{highlights}
\end{projectentry}

\vspace{0.4 cm}

\begin{projectentry}{Plug-n-Play Sensor Fusion Framework (P2SIF) for Autonomous Driving}{Jun 2024 -- Present}{With Prof. Dinesh Bharadia at UCSD}
    \begin{highlights}
        \item Developing a modular framework for fusing Radar, LiDAR, and Camera via text-based specialist models for LLM-driven decision making.
        \item Designing knowledge fusion mechanisms that decouple domain expertise (sensor processing) from generalist reasoning, to improve robustness and interpretability in autonomous driving.
        \item Currently evaluating system performance under adverse weather, plug-n-play sensor addition, and error injection cases, with focus on reliability and scalability.
        \item \textbf{Skills:} Modular sensor fusion, multi-sensor integration (Radar/LiDAR/Camera), LLM-driven reasoning, robustness and interpretability in autonomous driving.
    \end{highlights}
\end{projectentry}

\vspace{0.4 cm}

\begin{projectentry}{Realistic Radar Simulation Framework for CARLA (C-Shenron)}{Aug 2024 -- Feb 2025}{With Prof. Dinesh Bharadia at UCSD}
    \begin{highlights}
        \item Developed a high-fidelity radar simulation framework integrated into CARLA simulator for autonomous driving.
        \item This simulator incorporated material-aware reflectivity and configurable radar parameters.
        \item Built scalable data collection pipelines (850K+ frames across diverse towns, weather, and sensor setups) using Kubernetes for large-scale parallel simulation. 
        \item Trained and evaluated multi-radar fusion strategies along with ablations—demonstrating radar–camera fusion achieves performance comparable to LiDAR–camera baselines across CARLA leaderboard metrics.
        \item Published results in ACM SenSys’25 (Demo) and IEEE VTC’25 (Full Paper).
        % \item Skills: CARLA simulation, Docker, Kubernetes automation, large-scale data collection and training.
    \end{highlights}
\end{projectentry}

\vspace{0.4 cm}

\begin{projectentry}{Graph Learning and Data Inpainting via Deep Neural Networks}{Sept 2022 -- Jan 2024}{With Prof. Aditya Siripuram at IIT Hyderabad}
    \begin{highlights}
        \item Devised an algorithm for simultaneous data restoration and constructing a pertinent graph structure.
        \item Implemented a closed-loop feedback mechanism for graph learning guided by the inpainting performance.
        \item Achieved highest F1-Score of 0.92 (learned and actual graph) among existing methods on synthetic data. Demonstrated robust performance under 10dB noise and above on temperature sensor network datasets.
        \item Published in IEEE Signal Processing Letters (\href{https://doi.org/10.1109/LSP.2024.3501273}{10.1109/LSP.2024.3501273}) and presented at ICASSP 2025.
    \end{highlights}
\end{projectentry}

\vspace{0.4 cm}

\begin{projectentry}{New Training Algorithm for Deep Neural Networks}{Jan 2023 -- Jun 2024}{With Prof. Ayon Borthakur at IIT Hyderabad}
    \begin{highlights}
        \item Proposed a novel backpropagation-free training algorithm using a single forward pass with local loss functions. 
        \item Eliminates gradient propagation and activation storage, yielding significant memory and computational savings.
        \item Implemented on Multilayer MLPs, CNNs and LSTMs while demonstrating competitive performance with backpropagation on benchmarks such as MNIST, CIFAR-10, CIFAR-100, and SVHN.  
        \item Published in Transactions on Machine Learning Research (TMLR)
    \end{highlights}
\end{projectentry}\

\section{Selected Publications}

\publicationEntry
  {A Realistic Radar Simulator for End-to-End Autonomous Driving in CARLA} % Title
  {https://wcsng.ucsd.edu/c-shenron/} % Link
  {
    \equalAuthor{Satyam Srivastava,} 
    \equalAuthor{Jerry Li,} 
    \equalAuthor{\textbf{\textit{Pushkal Mishra}},} 
    Kshitiz Bansal, 
    Dinesh Bharadia
  } % Authors
  {IEEE 102nd Vehicular Technology Conference} % Conference
  {Oct 2025} % Date
  {Chengdu, China} % Location

\publicationEntry
  {Demo: Radar Simulator for CARLA} % Title
  {https://doi.org/10.1145/3715014.3724379} % Link
  {
    \textbf{\textit{Pushkal Mishra}}, 
    Satyam Srivastava, 
    Jerry Li, 
    Kshitiz Bansal, 
    Dinesh Bharadia
  } % Authors
  {23rd ACM Conference on Embedded Networked Sensor Systems} % Conference
  {May 2025} % Date
  {Irvine, USA} % Location

\publicationEntry
  {Learning Using a Single Forward Pass} % Title
  {https://openreview.net/forum?id=EDQ8QOGqjr/} % Link
  {
    \equalAuthor{Aditya Somasundaram}, 
    \equalAuthor{\textbf{\textit{Pushkal Mishra}}}, 
    Ayon Borthakur
  } % Authors
  {Transactions on Machine Learning Research} % Conference
  {Jun 2025} % Date
  {} % Location

\publicationEntry
  {Inpainting-Driven Graph Learning via Explainable Neural Networks} % Title
  {https://ieeexplore.ieee.org/abstract/document/10756724} % Link
  {
    \equalAuthor{Subbareddy Batreddy}, 
    \equalAuthor{\textbf{\textit{Pushkal Mishra}}}, 
    Yaswanth Kakarla, 
    Aditya Siripuram
  } % Authors
  {IEEE Signal Processing Letters, Presented at IEEE ICASSP 2025} % Conference
  {Apr 2025} % Date
  {Hyderabad, India} % Location

\vspace{-0.2cm}
{\footnotesize \raggedleft \textsuperscript{*} Equal Contribution.\par}
\vspace{-0.2cm}

\section{Awards}

\awardEntry
    {Electrical and Computer Engineering Department Fellowship}
    {Awarded full tuition support and stipend for five years of PhD at UC San Diego.}
    {2024}

\awardEntry
    {IEEE Signal Processing Cup}
    {Achieved rank 30 in the IEEE Signal Processing Cup challenge held worldwide.}
    {2023}

\awardEntry
    {Andy Grove Scholarship Award}
    {Recipient of the scholarship from Intel among 1800 competitive applicants to support undergraduate study.}
    {2022}

\awardEntry
    {All India Rank of 1910 in JEE Advanced}
    {JEE Advanced examination among 1 Million candidates}
    {2020}

\awardEntry
    {Merit Certificate in Physics (CBSE)}
    {Scored within top 0.1\% nationwide among 1.7 Million candidates}
    {2020}

\section{Skills}

\skillrow{Programming Languages}{Python, C/C++, MATLAB, Bash, Zsh, LaTeX, HTML, Verilog.}

\skillrow{Tools \& Frameworks}{PyTorch, TensorFlow, Docker, Git, ROS2, MATLAB DSP/Deep Learning Toolbox, Cadence, LTSpice, NGSpice, KiCAD.}

\skillrow{Systems Skills}{Low-level systems programming, Scripting Automation, Kubernetes, Embedded Programming (ESP32, Jetson Nano), CARLA Simulation framework.}


\section{Work Experience}

\experienceEntry
  {Signal Processing Intern}
  {Texas Instruments}
  {Bangalore}
  {May 2023 -- July 2023}
  {\begin{highlights}
      \item Designed and implemented a cross-talk cancellation algorithm for spatial audio playback.
      \item Developed an efficient low-order filter approximation using bi-quad and all-pass filters within a 3 dB tolerance.
      \item Proposed a new cross-talk cancellation architecture effective across different listener angles.
      \item Introduced performance metrics for evaluating spatial image reconstruction and cross-talk cancellation efficiency.
      \item Fully implemented solutions in MATLAB's Signal Processing and Deep Learning toolbox.
  \end{highlights}}

\experienceEntry
  {Teaching Assistant}
  {With Prof. Aditya Siripuram}
  {IIT Hyderabad}
  {Aug 2023 -- May 2024}
  {\begin{highlights}
      \item Assistant for both Linear Systems \& Signal Processing and Probability for AI.
      \item Conducted tutorial sessions for students to solve assignment and exam doubts.
      \item Formulated questions for bi-weekly quizzes and monthly exams.
  \end{highlights}}


\experienceEntry
  {Senior Club Member}
  {Epoch(ML) and Robotix Club}
  {IIT Hyderabad}
  {June 2021 -- May 2023}
  {\begin{highlights}
      \item Built a mini-bipedal robot from scratch using 3D-printed parts, servo motors, and Arduino drivers.
      \item Developed hand gesture recognition using motion sensors and DNN for time-series action classification.
      \item Conducted workshops and hands-on sessions on Arduino and motor drivers.
      \item Organized hackathons and knowledge-sharing sessions for the IIT Hyderabad community.
      \item Participated in research paper presentation sessions.
  \end{highlights}}
  

\experienceEntry
  {Class Representative}
  {EE Department, Undergraduate 2020 Batch}
  {IIT Hyderabad}
  {Aug 2022 -- May 2023}
  {\begin{highlights}
      \item Elected representative for the undergraduate batch of 2020.
  \end{highlights}}


\section{Selected Course Projects}

\begin{projectentry}{Autonomous Driving with Ultra-Wide Band (UWB) and Aruco Marker}{Apr 2025 -- Jun 2025}{ECE-148: Intro to Autonomous Driving Project at UCSD (Team of 4)}
    \begin{highlights}
        \item Designed and implemented an autonomous car-following system integrating real-time UWB-based triangulation, IMU fusion, and Aruco marker vision for robust localization and tracking even under non-line-of-sight conditions.
        \item Built ROS2 perception with: OAK-D Lite for marker detection, UWB based anchor triangulation, and Aruco-based following with curved-path obstacle navigation. 
        \item Project webpage: \href{https://github.com/UCSD-ECEMAE-148/148-spring-2025-final-project-team-10/}{https://github.com/UCSD-ECEMAE-148/148-spring-2025-final-project-team-10/}
        % \item Skills: Embedded systems, ROS2, Jetson Nano, ESP32 programming, real-time sensor fusion, Python/C++, Aruco vision, UWB localization.
    \end{highlights}
\end{projectentry}



\begin{projectentry}{IEEE 802.11 PHY/MAC Implementation}{Nov 2024 -- Jan 2025}{Course Projects in ECE 257A and ECE 257B at UCSD}
    \begin{highlights}
        \item Implemented core PHY components of IEEE 802.11 standard, including modulation/demodulation, OFDM channel access, channel estimation, and packet management (preamble detection, synchronization, and data recovery).  
        \item Designed and tested end-to-end packet construction and decoding pipeline, covering STF/LTF generation, Tx–Rx frequency/time synchronization, and CSI-based channel equalization.  
        \item Extended to advanced OFDM receiver implementation: packet detection via auto/cross-correlation, CFO/SFO estimation \& correction, pilot-based phase error mitigation, and BER benchmarking under QPSK and 16QAM.  
        \item Evaluated robustness across different synchronization and channel estimation strategies, analyzing trade-offs in detection accuracy and bit error rate.
    \end{highlights}
\end{projectentry}



\begin{projectentry}{Anchor-based Wi-Fi CSI Localization}{Aug 2024 -- Dec 2024}{ECE-257A: Modern Communication Networks Project at UCSD (Team of 3)}
    \begin{highlights}
        \item Proposed a novel Anchor-based CNN architecture, incorporating anchor-assisted features to achieve sub-meter accuracy in dynamic multipath environments. 
        \item Built an end-to-end data collection and preprocessing pipeline using Turtlebot, Raspberry Pi anchors, and SLAM for ground-truth alignment; performed CSI phase/magnitude calibration and ToF correction.
        \item Collected and trained on vast amounts of real‑scenario data in dynamic environments.
        \item Demonstrated that DL with CSI significantly outperforms RSSI and classical approaches for indoor localization.
    \end{highlights}
\end{projectentry}


\end{document}